\chapter*{Preface}
\addcontentsline{toc}{chapter}{Preface}

% ---------------------------------------------------------------
% Preface. Written by the author. Springer monograph convention is
% to cover: who the book is for, what it contains, how to read it,
% notational/typographical conventions, and a brief acknowledgement.
% Keep each section short (one paragraph is plenty); the Ishtar AI
% voice is conversational and concrete.
% ---------------------------------------------------------------

\section*{Who this book is for}

[AUDIENCE PARAGRAPH. Name the readers explicitly: engineers moving
from classical MLOps to LLMs, platform teams starting to productionize
RAG and agent systems, technical leads deciding when a finetune is
worth the cost, practitioners who need to operate models under SLOs
and budgets. Explain one prerequisite assumption: comfortable with
production software, some familiarity with machine-learning training
lifecycles; deep NLP expertise is not required.]

\section*{How to read this book}

[READING-PATH PARAGRAPH. Say that Part I (Foundations) is best read
in order; Parts II--IV can be sampled by topic. Recommend that readers
new to LLMOps start with Chapters 1--3; readers with MLOps experience
can jump to Chapter 4. Mention the end-to-end Ishtar AI case study
that runs across chapters and is concluded in Chapter 12.]

\section*{What changed since conception}

[CONTEXT PARAGRAPH. Briefly note that LLMOps as a practice barely
existed when this project began and the field moved visibly during
writing; the book's stance prefers durable operational patterns over
tool-of-the-month specifics. Invite readers to treat every vendor
name as an example of a category, not an endorsement.]

\section*{Conventions}

Throughout this book:

\begin{itemize}
  \item \textbf{Ishtar AI vignettes} appear in gray boxes and describe
        concrete operational situations drawn from a fictional
        production system used as the running case study.
  \item \textbf{Best Practice}, \textbf{Pitfall}, \textbf{Definition},
        and \textbf{Checklist} boxes highlight one crisp,
        actionable idea.
  \item \textbf{Algorithms} describe operational procedures (e.g.,
        rollout gates, routing policies); \textbf{Listings} show
        configuration and code.
  \item Cross-references to figures, tables, algorithms, and listings
        use the abbreviated Springer forms: Fig., Table, Alg.,
        Listing.
\end{itemize}

\section*{Acknowledgements}

[ACKNOWLEDGEMENTS PARAGRAPH. See also the dedicated Acknowledgements
chapter that follows, which contains the full human acknowledgements
and the required Springer AI-use declaration. This section can be a
one-sentence pointer: ``Acknowledgements, including the declaration
of AI-assisted content required by Springer Nature, appear in a
dedicated section after this Preface.'']
